\chapter{{\fontsize{14pt}{21pt}\selectfont\MakeUppercase{АНАЛИЗ ПРЕДМЕТНОЙ ОБЛАСТИ И ОБОСНОВАНИЕ ВЫБОРА ТЕХНОЛОГИЙ}}}
\label{cha:analysis}

% ============================================================
%  ГЛАВА 1. АНАЛИЗ ПРЕДМЕТНОЙ ОБЛАСТИ И ОБОСНОВАНИЕ ВЫБОРА ТЕХНОЛОГИЙ
% ============================================================

\section{Исследование инфраструктуры нагрузочного тестирования и роли имитационных сервисов}

В современной практике обеспечения качества сложных программных комплексов нагрузочное тестирование занимает критически важное место~\cite{kulikov-testing, gost-56920-2024}. Оно позволяет гарантировать стабильность и управляемость информационной системы при высоких эксплуатационных нагрузках. Одной из ключевых проблем при организации таких испытаний является необходимость взаимодействия объекта тестирования с множеством внешних систем и сервисов.

Использование реальных внешних интеграций в среде тестирования часто оказывается невозможным или нецелесообразным по ряду причин:

\begin{itemize}
\item ограниченность ресурсов: внешние системы могут иметь жёсткие лимиты на количество запросов или требовать значительных финансовых затрат на эксплуатацию;
\item изоляция среды: для получения достоверных результатов необходимо исключить влияние задержек и сбоев сторонних систем на показатели исследуемого объекта;
\item риски изменения данных: проведение тестов на реальных интеграциях может привести к нежелательному изменению или повреждению информации в продуктовых базах данных.
\end{itemize}

Для решения этих проблем применяются имитационные сервисы («заглушки») -- специализированное программное обеспечение, которое воспроизводит поведение реальных компонентов, отвечая на запросы тестируемой системы заранее определённым образом~\cite{newman-microservices}. В контексте данной работы в качестве имитационных сервисов рассматриваются Java-приложения в форматах \texttt{.jar} и \texttt{.war}.

Роль имитационных сервисов в инфраструктуре нагрузочного тестирования заключается в следующем:

\begin{itemize}
\item обеспечение автономности стенда: заглушки позволяют полностью изолировать тестируемую систему от внешнего окружения, создавая предсказуемую среду;
\item эмуляция различных сценариев: имитационные сервисы могут воспроизводить не только корректные ответы, но и ошибки (HTTP 5xx, тайм-ауты), что необходимо для проверки механизмов отказоустойчивости;
\item контроль интенсивности трафика: использование механизмов ограничения запросов (Rate Limiting) позволяет моделировать поведение внешних систем при достижении ими предельных порогов производительности.
\end{itemize}

Однако при увеличении масштаба информационной системы количество необходимых заглушек возрастает, что порождает проблему управления распределённой инфраструктурой. Процессы развёртывания, контроля состояния (активна/остановлена) и маршрутизации трафика к десяткам имитационных сервисов на различных серверах в ручном режиме существенно замедляют подготовку к испытаниям.

Таким образом, для эффективного функционирования стендов нагрузочного тестирования необходима специализированная система. Она должна централизованно управлять жизненным циклом заглушек, обеспечивать их мониторинг и динамическое проксирование запросов, что особенно актуально в условиях перехода на отечественные операционные системы, такие как ОС Astra Linux.

\section{Анализ существующих подходов к управлению сервисами-заглушками и проблема импортозамещения}

Традиционным подходом к управлению имитационными сервисами (заглушками), разработанными на языке Java, является их развёртывание в специализированных контейнерах серверов приложений. В современной практике разработки и тестирования сложился ряд стандартных решений, среди которых наиболее распространёнными являются:

\begin{itemize}
\item Apache Tomcat -- наиболее лёгкий и популярный сервлет-контейнер, используемый для запуска веб-приложений. Он обеспечивает базовую поддержку спецификаций Java Servlet и JavaServer Pages;
\item WildFly (ранее JBoss Application Server) -- более мощная и полнофункциональная платформа, поддерживающая широкий спектр стандартов Java EE (Jakarta EE). Данное решение предоставляет расширенные возможности управления ресурсами и кластеризации~\cite{wildfly};
\item Eclipse GlassFish -- эталонная реализация платформы Java EE, предлагающая развитые инструменты администрирования через графическую консоль и командную строку~\cite{glassfish}.
\end{itemize}

Анализ данных решений позволяет выделить общие тенденции и характеристики классических серверов приложений:

\begin{itemize}
\item универсальность и полнофункциональность: все перечисленные платформы спроектированы для запуска сложных корпоративных информационных систем. Однако для узкоспециализированной задачи -- работы лёгковесных имитационных заглушек -- их функциональность является избыточной;
\item значительные накладные расходы: общим недостатком является высокое потребление оперативной памяти и ресурсов центрального процессора самой платформой. При необходимости запуска десятков или сотен заглушек на одном или нескольких хостах суммарные затраты ресурсов на поддержание среды выполнения становятся критическими;
\item сложность автоматизации и оркестрации: управление большим парком разрозненных экземпляров серверов приложений требует внедрения дополнительных тяжеловесных систем автоматизации. Штатные средства управления (консоли администрирования) ориентированы на ручное конфигурирование, что затрудняет их использование в динамических сценариях нагрузочного тестирования;
\item монолитная архитектура управления: классические подходы подразумевают жёсткую привязку приложения к конкретной конфигурации сервера, что снижает гибкость при необходимости быстрого изменения параметров запуска (например, JVM-аргументов) для отдельных сервисов.
\end{itemize}

На текущий момент в российской ИТ-индустрии наблюдается отчётливая тенденция к отказу от использования зарубежных проприетарных и открытых решений в пользу отечественного программного обеспечения или собственных разработок. Это обусловлено политикой импортозамещения и необходимостью обеспечения технологического суверенитета.

Большинство западных систем управления инфраструктурой и серверов приложений (таких как продукты Oracle или Red Hat) сталкиваются с проблемами поддержки, обновления и совместимости с российскими операционными системами. В частности, эксплуатация имитационных стендов в защищённой среде ОС Astra Linux требует инструментов, которые нативно поддерживают механизмы управления этой операционной системы и не зависят от зарубежных репозиториев и лицензий~\cite{burenin-astra}.

Таким образом, анализ существующих подходов показывает, что использование традиционных серверов приложений типа Tomcat или WildFly для управления распределённой сетью заглушек становится неэффективным. Современные требования диктуют необходимость перехода к более лёгким, гибким и отечественным решениям, способным обеспечить централизованное управление жизненным циклом сервисов в рамках распределённых кластеров.

\section{Сравнительный анализ инструментальных средств разработки и обоснование технологического стека}

Выбор программных средств для реализации системы управления имитационными сервисами является фундаментальным этапом проектирования. От корректности этого выбора зависят эксплуатационные характеристики будущей системы: предельная пропускная способность, латентность (время отклика), масштабируемость и устойчивость к пиковым нагрузкам.

В ходе анализа предметной области и целей исследования было определено, что разрабатываемая система должна обеспечивать бесперебойную оркестрацию интенсивного потока запросов от нагрузочных станций. Это накладывает жёсткие ограничения на архитектуру: управляющий контур не должен становиться «узким местом» (bottleneck) при маршрутизации трафика. Требование минимизации накладных расходов исключает использование технологий, склонных к блокировке вычислительных ресурсов при ожидании ответа от внешних систем.

Для обоснования выбора проводится многокритериальный анализ по трём направлениям: язык программирования, веб-фреймворк и система управления базами данных.

\subsection{Обоснование выбора языка программирования серверной части}

Для разработки backend-составляющей системы рассматривались три языка программирования, наиболее востребованных в сфере высоконагруженных систем (HighLoad) и системной автоматизации.

Java -- объектно-ориентированный язык программирования, являющийся отраслевым стандартом для корпоративных систем и «родной» средой для запуска самих имитационных сервисов (\texttt{.jar}).

Преимущества: строгая статическая типизация, снижающая количество ошибок на этапе компиляции; развитая экосистема библиотек и инструментов профилирования; поддержка полноценной многопоточности (multithreading).

Недостатки: высокая ресурсоёмкость виртуальной машины (JVM), требующая значительного объёма оперативной памяти; длительное время «холодного старта», что критично для динамических сред тестирования; громоздкий синтаксис для простых задач взаимодействия с процессами операционной системы.

Go (Golang) -- компилируемый многопоточный язык программирования, разработанный Google, ориентированный на создание эффективных сетевых сервисов.

Преимущества: высочайшая производительность, сопоставимая с C++, благодаря компиляции в нативный код; эффективная модель конкурентности (goroutines), позволяющая обрабатывать тысячи соединений с минимальными накладными расходами; отсутствие внешних зависимостей (компиляция в единый бинарный файл).

Недостатки: отсутствие гибкости при работе с динамическими структурами данных (JSON произвольной структуры) из-за строгой типизации; специфическая обработка ошибок и отсутствие исключений, что увеличивает объём шаблонного кода (boilerplate); более высокий порог входа по сравнению с интерпретируемыми языками.

Python -- высокоуровневый интерпретируемый язык программирования общего назначения с динамической строгой типизацией.

Преимущества: наличие мощных инструментов асинхронного программирования (библиотека \texttt{asyncio}~\cite{python-asyncio}), позволяющих эффективно обрабатывать задачи, ограниченные скоростью ввода-вывода (I/O-bound), такие как сетевые запросы или ожидание ответа от базы данных~\cite{ramalho-python}; глубокая интеграция с Linux: модуль \texttt{subprocess} предоставляет удобный интерфейс для управления процессами, сигналами и потоками ввода-вывода; высокая скорость разработки благодаря лаконичному синтаксису.

Недостатки: производительность интерпретатора ниже, чем у компилируемых языков; наличие Global Interpreter Lock (GIL), ограничивающего выполнение потоков на одном ядре CPU (нивелируется использованием асинхронного подхода)~\cite{beazley-gil,python-asyncio,beazley-python-cookbook}.

Вывод: в качестве основного языка выбран Python. Возможность использования асинхронного программирования в сочетании с нативными средствами управления процессами Linux позволяет обеспечить необходимую производительность, сохраняя гибкость и скорость разработки.

\subsection{Сравнительный анализ веб-фреймворков}

Учитывая необходимость обеспечения высокой пропускной способности системы, выбор архитектуры веб-фреймворка является критическим. Рассматривались три кандидата: Django, Flask и FastAPI.

Django -- полнофункциональный синхронный фреймворк уровня Enterprise, следующий принципу «Batteries included» (всё включено).

Преимущества: наличие встроенной панели администрирования, ORM (Object-Relational Mapping) и системы аутентификации «из коробки»; развитая экосистема плагинов и высокий уровень безопасности; строгая структура проекта, облегчающая поддержку крупных монолитных приложений.

Недостатки: монолитная архитектура и избыточный функционал создают значительные накладные расходы (overhead) на каждый запрос; синхронная модель обработки запросов плохо подходит для высоконагруженных микросервисов, так как блокирует потоки при ожидании ввода-вывода; сложность отключения неиспользуемых компонентов для облегчения системы.

Flask -- лёгковесный микрофреймворк, классическое решение, основанное на спецификации WSGI (Web Server Gateway Interface)~\cite{greenberg-flask}.

Преимущества: минимализм и модульность: разработчик подключает только необходимые библиотеки; низкий порог входа и простота создания простых сервисов; гибкость в выборе архитектурных паттернов.

Недостатки: использование синхронной (блокирующей) модели ввода-вывода -- при высокой нагрузке это приводит к исчерпанию пула потоков сервера и росту задержек; отсутствие встроенных механизмов валидации данных (требует подключения сторонних расширений); снижение производительности при большом количестве одновременных соединений (проблема C10k)~\cite{kerrisk-linux}.

FastAPI~\cite{fastapi, lubanovich-fastapi} -- современный асинхронный фреймворк, построенный на спецификации ASGI (Asynchronous Server Gateway Interface)~\cite{asgi-spec}.

Преимущества: полная поддержка асинхронности (\texttt{async}/\texttt{await}) и событийного цикла, что позволяет обрабатывать тысячи запросов в секунду на одном процессе; высокая производительность благодаря использованию Starlette и uvloop, сопоставимая с Go и Node.js; встроенная быстрая валидация данных и сериализация на базе Pydantic~\cite{pydantic}.

Недостатки: относительно молодая экосистема по сравнению с Django и Flask; требует понимания принципов асинхронного программирования.

Вывод: выбран фреймворк FastAPI~\cite{fastapi, uvicorn}. Это единственное решение в экосистеме Python, способное обеспечить стабильную обработку интенсивного трафика с минимальными задержками благодаря асинхронной неблокирующей архитектуре.

\subsection{Обоснование выбора системы управления данными (СУБД)}

Для обеспечения высокой отзывчивости системы подсистема хранения данных должна обладать минимальным временем отклика на чтение конфигураций и запись статусов. Сравнивались конкретные распространённые решения.

PostgreSQL -- объектно-реляционная СУБД с открытым исходным кодом~\cite{morgunov-postgresql}.

Преимущества: гарантия целостности данных (ACID) и надёжность хранения; мощный язык запросов SQL и поддержка сложных выборок (JOIN); широкие возможности расширяемости (типы данных, индексы).

Недостатки: накладные расходы на дисковые операции ввода-вывода и журналирование (WAL) замедляют отклик; избыточность функционала для хранения простых пар «ключ-значение» (конфигурации заглушек); потребность в строгом проектировании схемы данных (Schema migrations).

MySQL -- реляционная СУБД, популярная база данных, часто используемая в веб-приложениях.

Преимущества: высокая скорость операций чтения в простых сценариях; широкая распространённость и простота администрирования; поддержка различных движков хранения (Storage Engines).

Недостатки: как и любая реляционная СУБД, зависит от скорости дисковой подсистемы; блокировки таблиц или строк могут снижать конкурентность записи при высокой нагрузке; жёсткость схемы данных.

MongoDB -- документоориентированная NoSQL СУБД, система, хранящая данные в формате BSON (бинарный JSON)~\cite{sadalage-nosql}.

Преимущества: гибкая схема данных (Schemaless), идеально подходящая для хранения JSON-конфигураций; горизонтальная масштабируемость (шардинг) из коробки.

Недостатки: потребляет значительно больше памяти и дискового пространства по сравнению с аналогами; в стандартной конфигурации уступает In-Memory решениям по скорости отклика; слабые гарантии транзакционности в старых версиях.

Redis -- In-Memory хранилище структур данных~\cite{redis-docs, macedo-redis}, высокопроизводительная СУБД, работающая полностью в оперативной памяти.

Преимущества: экстремально низкая латентность (время доступа менее 1 мс) благодаря отсутствию обращений к диску при чтении; простая модель данных «Ключ-Значение», полностью соответствующая задаче хранения состояний процессов; нативная поддержка механизма Time-To-Live (TTL) для временных данных.

Недостатки: зависимость объёма хранимых данных от объёма оперативной памяти; риск потери данных при аварийном отключении питания (решается настройкой персистентности).

Вывод: выбрана СУБД Redis~\cite{redis-docs}. Использование In-Memory архитектуры обеспечивает максимальное быстродействие. Проблема сохранности данных решается включением режима AOF (Append Only File)~\cite{redis-aof, kleppmann-ddia}, который асинхронно сохраняет операции на диск, гарантируя восстановление конфигураций после перезагрузки без ущерба для производительности.

\subsection{Итоговый выбор средств разработки}

На основании проведённого анализа и выявленных требований к быстродействию и надёжности утверждён следующий стек технологий для реализации выпускной квалификационной работы:

\begin{itemize}
\item язык программирования: Python (с использованием \texttt{asyncio});
\item веб-фреймворк: FastAPI (сервер приложений Uvicorn);
\item СУБД: Redis (режим персистентности AOF);
\item целевая платформа: ОС Astra Linux Special Edition~\cite{astra-linux}.
\end{itemize}

Данный выбор обеспечивает оптимальный баланс между производительностью, скоростью разработки и надёжностью функционирования системы в условиях высокой нагрузки.

\section{Обзор аналогичных решений и обоснование разработки собственного программного комплекса}

В рамках анализа предметной области необходимо рассмотреть существующие на рынке инструменты, применяемые для создания, эксплуатации и управления имитационными сервисами (заглушками). Поскольку разрабатываемая система предназначена для управления инфраструктурой заглушек в процессе нагрузочного тестирования, к аналогам следует отнести программные продукты, обеспечивающие виртуализацию сервисов (Service Virtualization), мокирование (Mocking), а также серверы приложений, традиционно используемые для хостинга Java-артефактов.

Сравнительный анализ проводится по критериям, критически важным для обеспечения стабильности нагрузочного стенда и автоматизации процессов тестирования:

\begin{itemize}
\item архитектура и потребление ресурсов: способность работать в условиях ограниченных аппаратных мощностей при запуске десятков экземпляров заглушек;
\item возможности управления процессами (Orchestration): наличие инструментов для удалённого запуска, остановки и обновления версий приложений-заглушек;
\item функциональность управления трафиком: наличие встроенных механизмов ограничения нагрузки (Rate Limiting) и эмуляции сетевых задержек;
\item совместимость и импортозамещение: возможность автономной работы в среде ОС Astra Linux без зависимости от зарубежных облачных сервисов и проприетарных лицензий.
\end{itemize}

\subsection{Анализ решения WireMock}

WireMock -- один из наиболее распространённых инструментов с открытым исходным кодом для виртуализации HTTP-сервисов. Продукт позволяет создавать заглушки, настраивать правила сопоставления запросов (request matching) и формировать динамические ответы.

Согласно официальной документации~\cite{wiremock}, WireMock поддерживает гибкую настройку ответов через JSON API, имитацию задержек (network latency) и ошибок (fault injection). Он может работать в двух режимах: как библиотека, встраиваемая в код автотестов, или как отдельный процесс (Standalone).

Недостатки в контексте поставленной задачи:

\begin{itemize}
\item ресурсоёмкость: WireMock разработан на Java и выполняется в виртуальной машине JVM. Запуск отдельного экземпляра WireMock для каждого имитируемого микросервиса создаёт значительные накладные расходы на оперативную память (RAM) и процессорное время (CPU). В условиях нагрузочного тестирования, когда требуется поднять 50--100 заглушек, суммарные затраты ресурсов на поддержание работы самих заглушек становятся неприемлемыми;
\item отсутствие встроенной оркестрации: в бесплатной версии (Open Source) WireMock отсутствует централизованная панель управления распределённым кластером. Инструмент фокусируется на логике ответа, а не на эксплуатации сервера;
\item ограничения Rate Limiting: базовый функционал позволяет имитировать задержки, однако полноценный алгоритм ограничения количества запросов в секунду (RPS) с отбрасыванием лишнего трафика требует написания кастомных расширений (extensions) на Java~\cite{wiremock}.
\end{itemize}

\subsection{Анализ решения MockServer}

MockServer -- популярное решение для мокирования систем, взаимодействующих по протоколам HTTP или HTTPS. Продукт часто используется в контейнеризированных средах (Docker, Kubernetes) для интеграционного тестирования.

MockServer предоставляет мощный API для создания «ожиданий» (Expectations) и верификации запросов. Согласно документации~\cite{mockserver}, он поддерживает проксирование трафика (Port Forwarding), запись проходящих запросов для последующего анализа и генерацию ответов на основе шаблонов (Mustache, Velocity).

Недостатки в контексте поставленной задачи:

\begin{itemize}
\item монолитная архитектура управления: MockServer хранит состояния «ожиданий» в памяти конкретного экземпляра приложения. Для создания распределённой системы требуется сложная настройка кластеризации, что утяжеляет инфраструктуру тестового стенда;
\item сложность динамического управления: MockServer ориентирован на программное управление через API во время прогона тестов. Использование его как постоянно действующего инфраструктурного компонента требует разработки дополнительной обвязки для мониторинга его состояния и автоматического перезапуска (Self-healing)~\cite{mockserver};
\item зависимость от Java: аналогично WireMock, продукт написан на Java, что влечёт проблемы с «холодным стартом» и высоким потреблением памяти (Heap Size), снижающим плотность размещения заглушек на сервере.
\end{itemize}

\subsection{Анализ решения Mountebank}

Mountebank -- инструмент с открытым исходным кодом, разработанный на платформе Node.js. Его ключевой особенностью является мультипротокольность: помимо HTTP/HTTPS, он поддерживает TCP, SMTP и другие протоколы через систему плагинов.

Mountebank использует концепцию «самозванцев» (imposters) -- лёгковесных сервисов, слушающих определённые порты. Документация~\cite{mountebank} указывает на возможность использования JavaScript-инъекций для реализации сложной логики обработки ответов.

Недостатки в контексте поставленной задачи:

\begin{itemize}
\item проблемы безопасности: возможность исполнения произвольного JS-кода (injection) в теле запроса создаёт уязвимости в контуре нагрузочного тестирования;
\item отсутствие управления артефактами: Mountebank позволяет создать логику ответа (скрипт), но не предназначен для запуска и управления бинарными файлами (заглушками в виде \texttt{.jar}), предоставленными разработчиками;
\item производительность в однопоточном режиме: в стандартной конфигурации Mountebank работает в одном потоке (Event Loop). При высокой нагрузке (тысячи RPS), характерной для нагрузочного тестирования, сложные вычисления внутри одной заглушки могут заблокировать обработку запросов других заглушек, работающих в том же экземпляре Mountebank~\cite{mountebank}.
\end{itemize}

\subsection{Анализ решения Apache Tomcat}

Apache Tomcat -- свободно распространяемый контейнер сервлетов, реализующий спецификации Jakarta Servlet и Jakarta Pages. Это классическое решение для хостинга Java-приложений, широко применяемое в промышленной эксплуатации.

Согласно официальной документации~\cite{tomcat}, Tomcat предоставляет развитые средства управления жизненным циклом веб-приложений. Компонент Manager App позволяет развёртывать (deploy), останавливать и перезапускать приложения без остановки сервера. Поддерживается кластеризация для репликации сессий~\cite{tomcat-cluster}. Для защиты от перегрузок предусмотрен фильтр RateLimitFilter, позволяющий ограничивать количество запросов с одного IP-адреса~\cite{tomcat-ratelimit}.

Недостатки в контексте поставленной задачи:

\begin{itemize}
\item избыточность архитектуры (Overhead): Tomcat является полноценным сервлет-контейнером. Использование его для запуска примитивных заглушек нерационально с точки зрения потребления ресурсов. Запуск отдельного экземпляра Tomcat для каждой заглушки потребляет сотни мегабайт оперативной памяти только на нужды самого сервера, что критично при дефиците ресурсов;
\item требования к формату артефактов: Tomcat работает с WAR-архивами. Однако современные микросервисы (и заглушки для них) чаще всего поставляются в виде Fat JAR (со встроенным сервером, например, Netty или Undertow). Для запуска таких заглушек в Tomcat требуется их пересборка и модификация кода, что усложняет процесс;
\item статичность конфигурации ограничений: механизмы Rate Limiting в Tomcat конфигурируются декларативно через XML-файлы при старте. В задачах нагрузочного тестирования требуется динамическое изменение параметров дросселирования трафика «на лету» (runtime) без перезагрузки сервера, что в Tomcat не реализовано «из коробки».
\end{itemize}

\subsection{Сравнительная характеристика}

Для наглядного сравнения рассмотренных решений с проектируемой системой составлена таблица~\ref{tab:comparison}.



% \begin{table}[H]
% \caption{Сравнительный анализ средств управления имитационными сервисами}
% \label{tab:comparison}
% \centering
% \small % Уменьшаем размер шрифта для таблицы
% \begin{tblr}{
%   colspec = {|X[1.2,l]|X[0.8,l]|X[0.8,l]|X[0.8,l]|X[0.9,l]|X[1.3,l]|},
%   hlines,
%   width = \linewidth
% }
% \textbf{Критерий} & \textbf{Wire\-Mock} & \textbf{Mock\-Server} & \textbf{Mounte\-bank} & \textbf{Apache Tomcat} & \textbf{Разрабаты\-ваемая система} \\
% Управление JAR/WAR без пересборки & Нет & Нет & Нет & Частично & Да \\
% Централи\-зованная оркестрация кластера & Нет & Нет & Нет & Нет & Да \\
% Динамический Rate Limiting & Нет & Нет & Нет & Нет & Да \\
% Потребление ресурсов управляющего слоя & Высокое (JVM) & Высокое (JVM) & Среднее & Высокое (JVM) & Низкое (Python) \\
% Совместимость с Astra Linux & Частичная & Частичная & Да & Частичная & Да \\
% \end{tblr}

% \noindent Источники данных:~\cite{wiremock},~\cite{mockserver},~\cite{mountebank},~\cite{tomcat}.
% \end{table}

% \begin{longtblr}[
%   caption = {Сравнительный анализ средств управления имитационными сервисами},
%   label   = {tab:comparison},
%   note{}  = {Источники данных:~\cite{wiremock},~\cite{mockserver},~\cite{mountebank},~\cite{tomcat}.},
% ]{
%   colspec  = {|X[1.4,l]|X[0.9,l]|X[0.9,l]|X[0.9,l]|X[1.0,l]|X[1.3,l]|},
%   hlines,
%   width    = \linewidth,
%   rowhead  = 1,
%   row{1}   = {font=\bfseries},
%   row{2-Z} = {font=\small},
%   row{1}   = {font=\bfseries\small},
% }
% Критерий & Wire\-Mock & Mock\-Server & Mounte\-bank & Apache Tomcat & Разрабаты\-ваемая система \\
% Управление JAR/WAR без пересборки & Нет & Нет & Нет & Частично & Да \\
% Централи\-зованная оркестрация кластера & Нет & Нет & Нет & Нет & Да \\
% Динамический Rate Limiting & Нет & Нет & Нет & Нет & Да \\
% Потребление ресурсов управляющего слоя & Высокое (JVM) & Высокое (JVM) & Среднее & Высокое (JVM) & Низкое (Python) \\
% Совместимость с Astra Linux & Частичная & Частичная & Да & Частичная & Да \\
% \end{longtblr}

% \begin{xltabular}{\textwidth}{|X|X|X|X|X|X|}
% \caption{Сравнительный анализ средств управления имитационными сервисами} \label{tab:comparison} \\
% \hline
% \textbf{Критерий} & \textbf{WireMock} & \textbf{MockServer} & \textbf{Mountebank} & \textbf{Apache Tomcat} & \textbf{Разраб. система} \\ \hline
% \endfirsthead
% \hline
% \textbf{Критерий} & \textbf{WireMock} & \textbf{MockServer} & \textbf{Mountebank} & \textbf{Apache Tomcat} & \textbf{Разраб. система} \\ \hline
% \endhead

% Управление JAR/WAR без пересборки & Нет & Нет & Нет & Частично & Да \\ \hline
% Централизованная оркестрация & Нет & Нет & Нет & Нет & Да \\ \hline
% Динамический Rate Limiting & Нет & Нет & Нет & Нет & Да \\ \hline
% Потребление ресурсов & Высокое (JVM) & Высокое (JVM) & Среднее & Высокое (JVM) & Низкое \\ \hline
% Совместимость с Astra Linux & Частично & Частично & Да & Частично & Да \\ \hline
% \end{xltabular}

{\small
\begin{xltabular}{\textwidth}{|p{3.2cm}|X|X|X|X|X|}
\caption{Сравнительный анализ средств управления имитационными сервисами} \label{tab:comparison} \\
\hline
\textbf{Критерий} & \textbf{Wire\-Mock} & \textbf{Mock\-Server} & \textbf{Mounte\-bank} & \textbf{Apache Tomcat} & \textbf{Разраб. система} \\ \hline
\endfirsthead
\hline
\textbf{Критерий} & \textbf{Wire\-Mock} & \textbf{Mock\-Server} & \textbf{Mounte\-bank} & \textbf{Apache Tomcat} & \textbf{Разраб. система} \\ \hline
\endhead

Управление JAR/WAR без пересборки & Нет & Нет & Нет & Частично & Да \\ \hline
Централизо\-ванная оркестрация & Нет & Нет & Нет & Нет & Да \\ \hline
Динами\-ческий Rate Limiting & Нет & Нет & Нет & Нет & Да \\ \hline
Потребление ресурсов & Высокое (JVM) & Высокое (JVM) & Среднее & Высокое (JVM) & Низкое \\ \hline
Совмести\-мость с Astra Linux & Частично & Частично & Да & Частично & Да \\ \hline
\end{xltabular}
}

\subsection{Обоснование разработки собственного программного комплекса}

Проведённый анализ показывает, что существующие на рынке решения делятся на две категории:

\begin{itemize}
\item инструменты логической виртуализации (WireMock, MockServer, Mountebank) -- отлично справляются с задачей «что ответить на запрос», но не решают задачу управления инфраструктурой (запуск, перезапуск, обновление версий приложений);
\item серверы приложений (Apache Tomcat) -- умеют управлять жизненным циклом, но слишком тяжеловесны, требуют специфического формата упаковки (WAR) и сложны в динамической настройке сетевых параметров.
\end{itemize}

Ни одно из рассмотренных решений не предоставляет функционала «лёгковесного агента» для управления сторонними Java-процессами «как есть» (без перепаковки).

Разработка собственной автоматизированной информационной системы обоснована следующими факторами.

Специфика задачи (Lifecycle Management): разрабатываемая система решает задачу оркестрации процессов операционной системы. Она позволяет загружать бинарный файл заглушки (JAR), полученный от разработчиков, и управлять его исполнением на удалённом сервере. Это исключает трудозатраты на пересборку заглушек под WireMock или Tomcat.

Эффективность использования ресурсов: перенос логики управления на Python и использование асинхронного фреймворка FastAPI позволяет минимизировать накладные расходы самого управляющего контура. Система работает поверх заглушек, не внедряясь в их память, в отличие от тяжёлых Java-контейнеров.

Централизация и удобство эксплуатации: реализация архитектуры «Master-Agent» объединяет разрозненные серверы в единый кластер. Единая панель управления (Dashboard) позволяет инженеру тестирования массово обновлять версии заглушек и менять настройки лимитов (RPS) в реальном времени, что невозможно при использовании Standalone-решений.

Технологическая независимость и импортозамещение: система разрабатывается на базе открытых технологий, адаптирована для работы в ОС Astra Linux и не зависит от внешних репозиториев (Maven Central, Docker Hub) или лицензионных ключей, что гарантирует стабильность работы в закрытых корпоративных контурах РФ.

\subsection {Вывод по подразделу}

Разработка собственного программного комплекса является единственным целесообразным решением, закрывающим пробел между инструментами логического мокирования и средствами управления инфраструктурой, обеспечивая высокую производительность и гибкость конфигурации тестовых сред.

\section{Формирование функциональных и технических требований к системе}

На основании проведённого анализа предметной области и специфики задач нагрузочного тестирования сформированы требования к разрабатываемому программному комплексу~\cite{wiegers-requirements, gost-34-602-2020}. Система классифицируется как инструмент автоматизации инфраструктуры нагрузочного тестирования. Архитектура решения должна быть гибкой, допускающей работу как в распределённом контуре, так и в рамках одного физического или виртуального сервера.

\subsection{Требования к режимам функционирования}

Система должна поддерживать два сценария развёртывания без изменения кодовой базы:

\begin{itemize}
\item распределённый режим (Cluster Mode): классическая клиент-серверная архитектура, где управляющий узел (Master) администрирует множество удалённых узлов (Agents), на которых выполняются заглушки. Взаимодействие осуществляется по сети;
\item автономный режим (Standalone Mode): консолидированный запуск всех компонентов системы на одном сервере. В этом режиме управляющий модуль и модуль исполнения заглушек функционируют в пределах одной операционной системы, обеспечивая изоляцию контура тестирования без необходимости сетевого взаимодействия между узлами инфраструктуры.
\end{itemize}

\subsection{Функциональные требования}

Функциональные требования определяют перечень задач, которые система должна выполнять для обеспечения процессов эмуляции сервисов.

\textbf{Требования к подсистеме управления жизненным циклом (Runtime).}

\begin{itemize}
\item поддержка форматов артефактов: система должна обеспечивать запуск Java-приложений, поставляемых в форматах исполняемых файлов \texttt{.jar} и веб-архивов \texttt{.war};
\item управление процессами: обеспечение запуска, корректной остановки и принудительного завершения процессов заглушек;
\item конфигурация запуска: возможность задания аргументов виртуальной машины (JVM Options), включая параметры памяти (\texttt{-Xms}, \texttt{-Xmx}) и системные свойства, через графический интерфейс управляющего узла (Master) перед стартом заглушки;
\item мониторинг: отслеживание статуса процесса (PID) и доступности его сетевого порта;
\item работа с логами: сбор стандартных потоков вывода (\texttt{stdout}, \texttt{stderr}) запущенных процессов и их трансляция в интерфейс системы для отладки.
\end{itemize}

\textbf{Требования к подсистеме проксирования запросов (Proxy Engine).}

Система должна работать как прозрачный прокси-сервер (Transparent Proxy) на прикладном уровне, реализуя следующий алгоритм:

\begin{enumerate}
\item приём входящего HTTP-запроса от внешней системы;
\item извлечение метаданных: HTTP-метод, заголовки, тело запроса и целевой путь (URI);
\item выполнение нового запроса к целевой заглушке (на её локальный порт) с использованием извлечённого метода, заголовков и тела;
\item получение ответа от заглушки;
\item возврат полученного ответа инициатору запроса в неизменном виде.
\end{enumerate}

Требования к управлению нагрузкой (Rate Limiting):

\begin{itemize}
\item наличие механизма ограничения количества запросов в секунду (RPS) для конкретной заглушки;
\item активация лимита производится через вызов конфигурационного эндпоинта системы (API);
\item при превышении лимита система должна возвращать HTTP-код 429 (Too Many Requests), не передавая запрос заглушке.
\end{itemize}

\textbf{Требования к управляющему модулю (Master/Dashboard).}

\begin{itemize}
\item управление артефактами: загрузка бинарных файлов (\texttt{.jar}, \texttt{.war}) в хранилище системы через веб-интерфейс;
\item реестр сервисов: отображение списка всех загруженных заглушек, их текущего статуса («Запущен», «Остановлен», «Ошибка») и адресов;
\item интерфейс конфигурации: предоставление полей ввода для настройки параметров запуска (JVM args) для каждого сервиса.
\end{itemize}

\subsection{Нефункциональные (технические) требования}

\textbf{Требования к производительности.}~\cite{gost-25010-2015}

\begin{itemize}
\item вносимая задержка (Overhead): накладные расходы на проксирование (время обработки запроса внутри системы без учёта времени работы самой заглушки) не должны превышать 20 мс;
\item пропускная способность: система должна обеспечивать высокую производительность проксирующего слоя. Снижение максимальной пропускной способности (RPS) заглушки при работе через систему не должно превышать 10\% относительно прямого обращения к заглушке.
\end{itemize}

\textbf{Требования к надёжности.}

\begin{itemize}
\item сохранение конфигурации и восстановление работоспособности: при перезапуске системы (или сервера) должны сохраняться настройки запуска заглушек и ссылки на загруженные артефакты (Persistence); для заглушек, находившихся в активном состоянии до перезапуска, система должна предпринять автоматическую попытку повторного запуска процесса;
\item изоляция сбоев: ошибка в одной заглушке не должна влиять на работоспособность остальных запущенных сервисов и самой управляющей системы.
\end{itemize}

\subsection{Требования к программно-аппаратному обеспечению}

\textbf{Системные требования.}

\begin{itemize}
\item операционная система: серверная часть должна функционировать под управлением ОС Astra Linux Special Edition (версия 1.7 и выше) или аналогичных Linux-дистрибутивов;
\item среда исполнения Java: поддержка OpenJDK 17~\cite{openjdk} для запуска пользовательских артефактов.
\end{itemize}

\textbf{Стек разработки.}

\begin{itemize}
\item язык программирования: Python (версия 3.10+) с использованием асинхронного подхода;
\item веб-фреймворк: FastAPI (для реализации высокопроизводительного ввода-вывода);
\item СУБД: Redis (для хранения конфигураций и реализации счётчиков Rate Limiter).
\end{itemize}

\subsection{Требования к информационной безопасности}

\begin{itemize}
\item контроль доступа: доступ к административному интерфейсу и API управления должен быть ограничен;
\item безопасность данных: загружаемые исполняемые файлы должны храниться в изолированных директориях с правами доступа, исключающими их несанкционированную модификацию третьими лицами.
\end{itemize}

\section {Выводы по первой главе}

В первой главе выпускной квалификационной работы был проведён комплексный анализ предметной области и обоснована необходимость создания специализированного инструментария для автоматизации нагрузочного тестирования.
В ходе исследования выявлено, что использование традиционных серверов приложений и существующих аналогов (WireMock, MockServer, Mountebank) для управления распределённой сетью заглушек сопряжено с избыточным потреблением ресурсов и сложностями в эксплуатации, а также не удовлетворяет требованиям по импортозамещению.

На основании сравнительного анализа обоснован выбор технологического стека: язык программирования Python с использованием асинхронного подхода, веб-фреймворк FastAPI для обеспечения высокой пропускной способности и In-Memory СУБД Redis для хранения состояний с минимальной задержкой.

Определена монолитная архитектура будущего решения с управлением удалёнными хостами через asyncssh/SFTP~\cite{fowler-eaa} и сформированы функциональные и технические требования к системе, включая поддержку форматов \texttt{.jar} и \texttt{.war}, механизмы ограничения нагрузки (Rate Limiting) и совместимость с ОС Astra Linux~\cite{gost-59793-2021}. Полученные результаты служат базисом для проектирования и технической реализации системы во второй главе.